3. Morality (Sila)

One who has faith in the Buddha, Dhamma, and Sangha and goes for Refuge in the Three Gems is said to be a Buddhist. The act of taking refuge is the pathway of entry into the Buddha's Dispensation or Teaching. However to become a virtuous lay Buddhist, one must refrain from immoral speech and actions by observing a moral code for the laity, called Five Precepts.

With regard to morality for the lay Buddhist, the Visuddhi Magga[3] or Path of Purification says: "Five Precepts as a permanent undertaking, Ten Precepts where possible and Eight Precepts on Uposatha Day, are the morality to be observed by the laity."

3.1 Five Precepts (Pancasila)

The Five Precepts form the actual practice of morality for the layman. They are the minimum ethical code, which are mandatory for all lay disciples. They are undertaken immediately after the taking of Refuge in the Buddha, Dhamma and Sangha at every Buddhist service or ceremony and are administered by a monk if one is present; otherwise the lay disciples can do it by themselves. It is usual for devout lay Buddhists to undertake the Five Precepts as part of their daily recitation. The Five Precepts are undertaken by reciting the following:

i) Panatipata veramani sikkhapadam samadiyami. (Taking of life, I undertake the training rule to abstain.)

ii) Adinnadana veramani sikkhapadam samadiyami. (Taking of things not given, I undertake the training rule to abstain.)

iii) Kamesu veramani sikkhapadam samadiyami. (Sexual misconduct, I undertake the training rule to abstain.)

iv) Musavada veramani sikkhapadam samadiyami. (False speech, I undertake the training rule to abstain.)

v) Sura meraya majja pamadatthana veramani sikkhapadam samadiyami. (Intoxicating liquors, spirits and drugs, I undertake the training rule to abstain.)

3.2 Dhamma Way to Compare Oneself with Another

Once when the Buddha was in the village of Veludvara[4] in Kosala country, he was asked by the villagers to teach them how to attain to the heavenly world where the virtuous are reborn. Thereupon the Buddha taught them to reflect on the Dhamma way to compare oneself with another, which leads to right understanding in the observance of the Five Precepts.

* On the matter of killing: Every person wishes to live and not to die; everyone is fond of pleasure and adverse to pain.

* If someone were to kill us, it would not be pleasing or delightful to us. Also if one kills another who wants to live and not to die, it would not be pleasing or delightful to the other person

* So something that is not pleasing to oneself must also not be pleasing to another. Therefore something that is not pleasing to one-self should not be imposed on another.

As a result of this reflection, he himself abstains from killing living beings. He encourages other so to abstain and he speaks in praise of so abstaining. Thus his bodily conduct is absolutely pure in these three aspects. By similar reflection and reasoning, one acquires a better understanding and appreciation of each of the precepts. In the case of the fifth precept, although the partaking of intoxicants appears to involve oneself only, it is the most dangerous because it can lead to the violation of all the precepts thereby causing more harm to onself and others.

Thus one who keeps the Five Precepts is an asset to oneself and others. In the Abhisanda Sutta of Anguttara Nikaya VIII, 39, the Buddha has compared the meticulous observance of the Five Precepts as five great gifts to oneself and others. By doing so, he gives freedom from fear, freedom from hostility, freedom from oppression to limitless numbers of beings. In giving freedom from fear, freedom from hostility, freedom from oppression to limitless numbers of beings, he gains a share in limitless from from fear, freedom from hostility, and freedom from oppression.

3.3 Consequences of Breaking & Keeping the Five Precepts
Laymen who are ignorant may consider the Five Precepts as something that prevent them from doing what they like, espescially the part on telling lies, slander or gossip and drinking beer or wine, which many non-Buddhists consider as socially acceptable and entertaining. Some even look upon the precepts as handcuffs to chain their hands and gags to shut their mouths, denying them the freedom to do what they like.

According to the Discourse on the Bad Effects of Evil Deeds in Anguttara Nikaya VIII, 40, the breaking of precepts when pursued, practised, increased, causes one to arise in hell, in the animal world and in the world of ghosts. If reborn as a human being, the following are the very least results:
* Killing will lead to shortening of one's life.
* Stealing will lead to loss of one's wealth.
* Sexual misconduct will breed rivalry and hatred.
* False speech will cause one to be falsely accused.
* Partaking of intoxicants will cause one to be afflicted with insanity.

On the other hand, the sincere meticulous observance of the Five Precepts leads to the accumulation of wholesome kamma tending to rebirth in the happy realms of humans or deities. If reborn as a human being, the following are the results:
* Abstentation from killing will lead to longevity.
* Abstentation from stealing will lead to prosperity.
* Abstentation from sexual misconduct will lead to popularity.
* Abstentation from lying will cause one to have a good reputation.
* Abstentation from partaking of intoxicants will lead to mindfulness and wisdom.

In the Mahaparinibanna Sutta[5], the Buddha expounded the villages of Pataligama the five blessings gained by one who observes the precepts and who is established in morality.
* He acquires much wealth as a result of his diligence.
* He acquires a good reputation and fame.
* He approaches and enters any assembly with complete self-confidence, without any fear or hesitation.
* He lives the full span of his life and dies undeluded.
* After death, he is reborn in the happy realms of humans or devas.

3.4 Protection through the Five Precepts

